\documentclass[12pt,a4paper]{article}
\usepackage[utf8]{inputenc}
\usepackage[T2A]{fontenc}
\usepackage[russian]{babel}
\usepackage{amsmath}
\usepackage{amsfonts}
\usepackage{amssymb}
\usepackage{graphicx}
\usepackage{geometry}
\usepackage{indentfirst}
\usepackage{multirow}
\usepackage{array}

\geometry{left=2.5cm, right=1.5cm, top=2cm, bottom=2cm}

\newcommand{\ReNum}{\mathrm{Re}}

\begin{document}

\begin{titlepage}
    \begin{center}
		{\large МОСКОВСКИЙ ФИЗИКО-ТЕХНИЧЕСКИЙ ИНСТИТУТ (НАЦИОНАЛЬНЫЙ ИССЛЕДОВАТЕЛЬСКИЙ УНИВЕРСИТЕТ)}
	\end{center}
	\begin{center}



        \vspace{6cm}
        {\huge
		\begin{center}
			{\bf Лабораторная работа 20}\\
             Гидродинамическая устойчивость вращательного течения Куэтта 
		\end{center}
	}
	
		
		\vspace{10cm}
	\begin{flushright}
	   {\large работу выполнили студенты группы Б03-305}\par
	   \vspace{0.3cm}
	   {\LARGE Кривушин Артём, Симухин Егор, Глебов Арсений}
    \end{flushright}
	   \vspace{2cm}
	   \begin{center}
		  Долгопрудный, 2026
	   \end{center}
    	   \end{center}
\end{titlepage}

\section{Цель работы}

Целью лабораторной работы является экспериментальное изучение гидродинамической устойчивости течения между двумя соосно вращающимися цилиндрами — устойчивости течения Куэтта.

Визуализация течения в отраженном свете осуществляется введением алюминиевой пудры в жидкость (конденсаторное масло).

В процессе выполнения работы студенты знакомятся с методами теоретического исследования гидродинамической устойчивости течений; описанием методики, измерений и аппаратуры; экспериментально определяют границу устойчивости однородного по оси $z$ вращательного течения при числах Рейнольдса до 4000; сопоставляют результаты с расчетными данными.


\section{Краткие сведения о цилиндрическом течении Куэтта}


При изучении течения в зазоре между двумя коаксиальными цилиндрами, вращающимися с разными угловыми скоростями (течение Куэтта), было установлено, что возможны два типа перехода от ламинарного к турбулентному режиму.

Первый имеет место, когда угловая скорость вращения внутреннего цилиндра больше, чем внешнего. Если скорость внешнего цилиндра зафиксировать, то при достижении внутренним цилиндром некоторой определенной скорости происходит потеря устойчивости однородного вдоль оси вращения стационарного течения; вместо него возникает устойчивое (при данной скорости), но периодическое по оси вращения течение — система вихрей в плоскости меридиана, имеющих чередующиеся противоположные вращения и занимающих все пространство между цилиндрами (рис. \ref{fig:20.1}).

Если цилиндры вращаются в противоположные стороны, образуется два ряда вихрей; из них один соприкасается с внешним, другой — с внутренним цилиндром (рис. \ref{fig:20.2}).

Увеличение скорости вращения приводит к потере устойчивости этого течения: возникают периодическая по кругу структура — волна, движущаяся вдоль "барашков" со скоростью, отличающейся от скорости вращения цилиндров (рис. \ref{fig:20.3}).

\begin{figure}[h]
	\centering
	\begin{minipage}[b]{0.3\textwidth}
		\centering
		\includegraphics[width=\textwidth]{20.1.png}
        		\caption{Вращающиеся цилиндры. Первый случай.}

	\end{minipage}
	\hfill
	\begin{minipage}[b]{0.3\textwidth}
		\centering
\includegraphics[width=\textwidth]{20.2.png}
		\caption{Вращающиеся цилиндры. Второй случай.}
\end{minipage}
	\hfill
	\begin{minipage}[b]{0.3\textwidth}
		\centering
\includegraphics[width=\textwidth]{20.3.png}	
		\caption{Вращающиеся цилиндры. Третий случай.}
\end{minipage}
\end{figure}

Более высокие моды, возникающие с дальнейшим увеличением скорости вращения, соответствуют гармоникам этих двух основных частот. Предполагается, что при достаточно больших числах Рейнольдса дискретный спектр частот переходит в непрерывный, что и соответствует переходу к турбулентному течению.

\section{Описание лабораторной установки}

Установка состоит из двух соосных цилиндров 1 и 2, приводимых во вращение электродвигателями 3 и 4, роторы которых установлены непосредственно на осях вращения цилиндров.

\begin{figure}[h!]
	\centering
	\begin{minipage}[b]{0.3\textwidth}
		\centering
		\includegraphics[width=\textwidth]{ust.png}
		\caption{Разрез установки}
		\label{fig:20.1}
	\end{minipage}
	\hfill
\end{figure}

Внешний цилиндр изготовлен из прозрачного материала (оргстекло).

Скорость вращения задается внешним звуковым генератором ГЗ-53. Сигнал генератора с некоторой частотой $f_2$ подается на управляющую схему в силовом блоке установки, обеспечивающую подачу на три группы обмоток статора соответствующего двигателя сдвинутые по фазе на $120^\circ$ импульсы тока от источника питания. Возникающее при этом вращающееся магнитное поле приводит во вращение ротор двигателя, связанный с соответствующим цилиндром. Изменение частоты сигнала генератора приводит к изменению скорости вращения двигателя. Связь эта не абсолютно жесткая: переменная нагрузка на двигатель вызывает различное проскальзывание ротора относительно магнитного поля. Это требует независимого от показаний частоты генератора измерения скорости вращения двигателя.

На роторе каждого двигателя установлен диск с N отверстиями, равномерно расположенными по кругу ($N = 100$). Свет от лампы накаливания, проходя через пробегающие мимо отверстия, вызывает импульсы тока в фоторезисторе, помещенном по другую сторону от диска. Эти импульсы подаются на частотомер ЧЗ-63, который рассчитывает число импульсов за секунду. Таким образом, если показание частотомера равно $f$, то частота вращения цилиндра равна:
\begin{equation}
	\Omega = \frac{f}{N} = \frac{f}{100}
\end{equation}

В установке предусмотрена возможность торможения внешнего цилиндра и вращения его в двух направлениях.



\section{Стационарное распределение скоростей в течении Куэтта}

Рассмотрим движение жидкости, заключенной между двумя коаксиальными бесконечными цилиндрами, вращающимися вокруг своей оси с угловыми скоростями $\Omega_1$ и $\Omega_2$; радиусы цилиндров пусть будут $R_1$ и $R_2$, причем $R_2 > R_1$. Выберем цилиндрические координаты $r, z, \varphi$ с осью $z$ по оси цилиндров. Из симметрии очевидно, что
\begin{equation}
v_z = v_r = 0, \quad v_\varphi = v(r); \quad p = p(r).
\end{equation}

Уравнение Навье-Стокса в цилиндрических координатах дает в рассматриваемом случае два уравнения:
\begin{align}
\frac{dp}{dr} &= \rho \frac{v^2}{r}, \label{eq:18.1} \\
\frac{d^2v}{dr^2} + \frac{1}{r} \frac{dv}{dr} - \frac{v}{r^2} &= 0. \label{eq:18.2}
\end{align}

Второе из этих уравнений имеет решения типа $r^n$; подстановка решения в таком виде дает $n = \pm 1$, так что
\begin{equation}
v = ar + \frac{b}{r}.
\end{equation}

Постоянные $a$ и $b$ находятся из предельных условий, согласно которым скорость жидкости на внутренней и внешней цилиндрических поверхностях должна быть равна скорости соответствующего цилиндра: $v = R_1 \Omega_1$ при $r = R_1$, $v = R_2 \Omega_2$ при $r = R_2$. В результате получаем распределение скоростей в виде
\begin{equation}
\boxed{v = \frac{\Omega_2 R_2^2 - \Omega_1 R_1^2}{R_2^2 - R_1^2} r + \frac{(\Omega_1 - \Omega_2) R_1^2 R_2^2}{R_2^2 - R_1^2} \frac{1}{r}.} \label{eq:18.3}
\end{equation}

Распределение давления получается отсюда согласно (\ref{eq:18.1}) простым интегрированием.

При $\Omega_1 = \Omega_2 = \Omega$ скорость $v = \Omega r$, т. е. жидкость вращается как целое вместе с цилиндрами. При отсутствии внешнего цилиндра ($\Omega_2 = 0; R_2 = \infty$) скорость
\begin{equation}
v = \frac{\Omega_1 R_1^2}{r}.
\end{equation}

\section{Ход работы}

Для начала запустим установку и дождемся равномерного распределения металлической пудры. Затем проведем наблюдения. 
Сначала с внешним заторможеным цилиндром, повышая частоту вращения, пронаблюдаем последовательное образование кольцевых структур в течении, запишем частоты необходимые для образования колец. Также можно заметить, что система обладает некоторым гистерезисом, поэтому можем собрать данные о появлении и исчезновении колец при увеличении и затем уменьшении частоты:

\begin{table}[h!]
    \centering
    \begin{tabular}{|c|c|c|c|c|c|c|c|c|c|c|c|}
    \hline
         $N$& 1 & 2 & 3 & 4 & 5 & 6\\ \hline
         $f, Гц$& 6.31 & 6.70 & 7.10 & 9.40 &9.55 & 9.63\\ \hline
    \end{tabular}
    \caption{При повышении частоты}
    \label{tab:placeholder}
\end{table}

\begin{table}[h!]
    \centering
    \begin{tabular}{|c|c|c|c|c|c|c|c|c|c|c|c|}
    \hline
         $N$& 1 & 2 & 3 & 4 & 5 & 6 & 7 & 12 &15 & 17\\ \hline
         $f, Гц$& 8.30 & 8.55 & 8.90 & 8.95 &9.00 & 9.15 & 9.30 & 9.30 &9.60 & 10.0\\ \hline
    \end{tabular}
    \caption{При понижении частоты}
    \label{tab:placeholder}
\end{table}

Данные отобразим на графиках:

\begin{figure}[h!]
    \centering
    \includegraphics[width=0.7\linewidth]{IMG_0027.png}
    \caption{Повышение частоты}
    \label{fig:placeholder}
\end{figure}

\begin{figure}[h!]
    \centering
    \includegraphics[width=0.7\linewidth]{IMG_0026.png}
    \caption{Понижение частоты}
    \label{fig:placeholder}
\end{figure}

\newpage

Также приведем внешний вид этой устойчивой формы течения:
\newpage

\begin{figure}[h!]
    \centering
    \includegraphics[width=0.7\linewidth]{IMG_0023.jpeg}
    \caption{Внешний цилиндр заторможен}
    \label{fig:placeholder}
\end{figure}

\newpage

Затем пустим второй цилиндр и регулируя частоты вращения внутреннего цилиндра найдем устойчивый режим при движении их обоих:

\begin{figure}[h!]
    \centering
    \includegraphics[width=0.4\linewidth]{IMG_0024.jpeg}
    \caption{в одну сторону}
    \label{fig:placeholder}
\end{figure}

Сделаем то-же самое при разнонаправленном вращении:

\begin{figure}[h!]
    \centering
    \includegraphics[width=0.4\linewidth]{IMG_0025.jpeg}
    \caption{В разные стороны}
    \label{fig:placeholder}
\end{figure}

\newpage

Также построим теоретическую кривую устойчивости, к сожалению мы не смогли снять достаточно данных чтобы построить ее экспериментально.

\begin{figure}[h!]
	\centering
	\begin{minipage}[b]{0.99\textwidth}
		\centering
		\includegraphics[width=\textwidth]{moe.png}
		\caption{Граница устойчивости вращательного течения Куэтта}
		\label{fig:20.1}
	\end{minipage}
	\hfill
\end{figure}

\section{Вывод}
В ходе данной работы было исследовано течение Куэтта


\end{document}